\noindent\textbf{Por qué se usa este procedimiento.} Bredt da flujo constante $q=T/(2A_m)$ en una célula cerrada. Como $A_m=\pi R^2$ y $\tau=q/t$, resulta $\tau=T/(2\pi R^2t)$. Las dimensiones son las de tensión.

\begin{enumerate}

\item El área encerrada es $A_m = \pi R^2$.

\end{enumerate}

\noindent\textbf{Operaciones y resultados intermedios.}

\begin{enumerate}

\item Tensión tangencial.

\[\fitmath{y_{1}=\frac{2000000}{2\cdot \pi\cdot {40}^{2}\cdot 1.5}\approx 132.6291\,\mathrm{MPa}}\]

\end{enumerate}
